% !TEX program = xelatex
% 贝塞尔函数：一般阶、递推、正交、圆膜
\documentclass[fontset=windows, chinese, 10pt, xcolor=dvipsnames]{beamer}
% Shared preamble for the chapter 11 standalone lectures.
% Each lecture sets \title and \subtitle after this file is read.

\usepackage{amsmath,amssymb,mathtools}
\usepackage{xeCJK}
\usepackage{zhnumber}
\usepackage{tikz}
\usepackage{tikz-3dplot}
\usetikzlibrary{arrows.meta,calc,positioning}

\usetheme{metropolis}
\usepackage{beamercolorthemeHistNanjing}

\newcommand{\red}[1]{\textcolor{red}{#1}}
\newcommand{\blue}[1]{\textcolor{blue}{#1}}
\newcommand{\green}[1]{\textcolor{green}{#1}}
\newcommand{\gray}[1]{\textcolor{gray}{#1}}
\newcommand{\orange}[1]{\textcolor{orange}{#1}}

\newcommand{\sssub}[1]{ {\mathrm{#1} }}  % script size subscript\\\
\newcommand{\funcd}[2]{ \frac{d{#1}}{d{#2}}} %function full derivative
\newcommand{\funcpd}[2]{ \frac{\partial{#1}}{\partial{#2}}} %function partial derivative
\newcommand{\fnald}[2] { \frac{\delta {#1}}{\delta {#2}}} % functional derivative

\renewcommand{\Re}{\operatorname{Re}}
\renewcommand{\Im}{\operatorname{Im}}
\def\ext{\sssub{ext}}
\def\x{\sssub{x}}
\def\xc{\sssub{xc}}
\def\H{\sssub{H}}
\def\Hx{\sssub{Hx}}
\def\Hxc{\sssub{Hxc}}
\def\c{\sssub{c}}
\def\s{\sssub{s}}

% \def\ext{\sssub{ext}}
% \def\ext{\sssub{ext}}
% \def\ext{\sssub{ext}}
% \def\ext{\sssub{ext}}
% \defsssub{ext}


% definitions of bold faced symbols
\def\brpppp{{\mathbf{r}^{\prime\prime\prime\prime}}}
\def\brppp{{\mathbf{r}^{\prime\prime\prime}}}
\def\brpp{{\mathbf{r}^{\prime\prime}}}
\def\brp{{\mathbf{r}^{\prime}}}
\def\bzp{{\mathbf{z}^{\prime}}}
\def\bxp{{\mathbf{x}^{\prime}}}
\def\tp{{{t}^{\prime}}}
\def\tpp{{{t}^{\prime\prime}}}
\def\tppp{{{t}^{\prime\prime\prime}}}

\def\tbr{{\tilde{\mathbf{r}}}}
\def\bk{{\mathbf{k}}}
\def\br{{\mathbf{r}}}
\def\bp{{\mathbf{p}}}
\def\bq{{\mathbf{q}}}
\def\bv{{\mathbf{v}}}
\def\bz{{\mathbf{z}}}
\def\mbf{{\mathbf{f}}}

\def\bx{{\mathbf{x}}}
\def\bR{{\mathbf{R}}}
\def\bM{{\mathbf{M}}}
\def\bP{{\mathbf{P}}}
\def\bT{{\mathbf{T}}}
\def\bK{{\mathbf{K}}}
\def\bA{{\mathbf{A}}}
\def\bB{{\mathbf{B}}}
\def\bD{{\mathbf{D}}}
\def\bE{{\mathbf{E}}}
\def\bF{{\mathbf{F}}}
\def\bG{{\mathbf{G}}}
\def\bH{{\mathbf{H}}}
\def\bI{{\mathbf{I}}}
\def\bJ{{\mathbf{J}}}
\def\bP{{\mathbf{P}}}
\def\bS{{\mathbf{S}}}
\def\bV{{\mathbf{V}}}
\def\bX{{\mathbf{X}}}
\def\bY{{\mathbf{Y}}}
\def\dd{{\mathop{}\!\mathrm{d}}}
\def\dr{{\!\!\dd\br\,}}
\def\dx{{\!\!\dd\bx\,}}
\def\dy{{\!\!\dd\by\,}}
\def\dk{{\!\!\dd\bk\,}}
\def\dG{{\!\!\dd\bG\,}}


\def\half{{\frac{1}{2}}}
\def\tr{{\mathrm{tr}}}


\def\Arg{{\mathrm{Arg}\;}}
\def\arg{{\mathrm{arg}\;}}
\def\Res{{\mathrm{Res}\;}}

% unit-vector basis (bold under XeLaTeX; \mathbf can look non-bold)
\newcommand{\ue}{\boldsymbol{e}}
\newcommand{\un}{\hat{\boldsymbol{n}}}


\author{罗凯}
\institute{南京理工大学物理学院}
\date{\zhdigits{\the\year},\zhganzhinian{\the\year}年秋}

\usebackgroundtemplate{%
\begin{tikzpicture}[remember picture,overlay]
  \node[opacity=0.05, at=(current page.center)] {
    \includegraphics[width=0.6\paperwidth]{njust-logo.jpeg}
  };
\end{tikzpicture}}


\title{特殊函数简介}
\subtitle{第三讲：贝塞尔函数}

\begin{document}

\begin{frame}
  \titlepage
\end{frame}

\begin{frame}{本讲提纲}
  圆膜、圆柱导热和圆波导的径向方程都是贝塞尔方程。整数阶的幂级数已经由弗罗贝尼乌斯方法得到，这里写成一般阶，并补上递推和正交性。
  \begin{enumerate}
    \item $J_{\nu}$ 的级数，整数阶时的第二类解；
    \item 生成函数和递推；
    \item 权重 $x$ 下的正交与模；
    \item 圆膜的本征频率，以及球贝塞尔函数。
  \end{enumerate}
\end{frame}

\section{级数}

\begin{frame}{方程}
  $\nu$ 阶贝塞尔方程是
  \begin{equation*}
    x^2 y''+x y'+(x^2-\nu^2)y=0.
  \end{equation*}
  $x=0$ 是正则奇点。第一类解
  \begin{equation*}
    \boxed{\displaystyle
      J_{\nu}(x)
      =\sum_{k=0}^{\infty}
        \frac{(-1)^k}{k!\,\Gamma(k+\nu+1)}
        \Bigl(\frac{x}{2}\Bigr)^{2k+\nu}.}
  \end{equation*}
  级数对所有 $x$ 收敛。$\nu=n$ 为非负整数时 $\Gamma(k+n+1)=(n+k)!$，回到前面的 $J_n$。
\end{frame}

\begin{frame}{第二个解}
  $\nu$ 不是整数时，$J_{\nu}$ 和 $J_{-\nu}$ 线性无关，通解是
  \[
    y=A J_{\nu}(x)+B J_{-\nu}(x).
  \]
  $\nu=n$ 为整数时 $J_{-n}=(-1)^n J_n$。第二个解是第二类贝塞尔函数 $Y_n$，在 $x=0$ 附近含 $\ln x$。包含原点的圆域里通常只留 $J_n$。
\end{frame}

\begin{frame}{虚宗量}
  把 $x$ 换成 $\mathrm{i}x$，方程里 $x^2$ 的符号反过来，
  \[
    x^2 y''+x y'-(x^2+\nu^2)y=0.
  \]
  原点有限的解是修正贝塞尔函数
  \[
    I_{\nu}(x)=\mathrm{i}^{-\nu}J_{\nu}(\mathrm{i}x).
  \]
  圆柱上的拉普拉斯方程，轴向取指数或双曲函数时，径向就是 $I_{\nu}$。
\end{frame}

\section{递推}

\begin{frame}{生成函数}
  整数阶有
  \begin{equation*}
    \exp\Bigl(\frac{x}{2}\Bigl(t-\frac{1}{t}\Bigr)\Bigr)
    =\sum_{n=-\infty}^{\infty} J_n(x)\,t^n,
    \qquad t\neq 0.
  \end{equation*}
  对 $x$ 求导，左边多一个因子 $\frac12(t-1/t)$。对 $t$ 求导，左边多一个因子 $(x/2)(1+1/t^2)$。比较 $t^n$ 的系数，
  \begin{align*}
    J_{n-1}(x)+J_{n+1}(x)&=\frac{2n}{x}J_n(x),\\
    J_{n-1}(x)-J_{n+1}(x)&=2J_n'(x).
  \end{align*}
\end{frame}

\begin{frame}{对一般阶也成立}
  把 $n$ 写成 $\nu$，两式仍成立，可以直接用级数核对。由此
  \begin{equation*}
    \frac{\mathrm{d}}{\mathrm{d}x}\bigl[x^{\nu}J_{\nu}(x)\bigr]
    =x^{\nu}J_{\nu-1}(x),
  \end{equation*}
  \[
    \frac{\mathrm{d}}{\mathrm{d}x}\bigl[x^{-\nu}J_{\nu}(x)\bigr]
    =-x^{-\nu}J_{\nu+1}(x).
  \]
  有限圆柱上对 $\rho J_0(k\rho)$ 积分，用的就是第一式在 $\nu=1$ 的积分形式 $\mathrm{d}[u J_1(u)]/\mathrm{d}u=u J_0(u)$。
\end{frame}

\begin{frame}{积分表示}
  令 $t=e^{\mathrm{i}\theta}$，生成函数变成
  \[
    e^{\mathrm{i}x\sin\theta}
    =\sum_{n=-\infty}^{\infty} J_n(x)\,e^{\mathrm{i}n\theta}.
  \]
  取实部，对 $\theta$ 从 $0$ 积到 $\pi$，
  \[
    J_n(x)
    =\frac{1}{\pi}\int_0^{\pi}
      \cos(n\theta-x\sin\theta)\,\mathrm{d}\theta.
  \]
  $x=0$、$n=0$ 时被积函数是 $1$，所以 $J_0(0)=1$。对 $x$ 在积分号下求导得 $J_0'=-J_1$，因而 $J_1(0)=0$，$J_1(x)\sim x/2$。
\end{frame}

\begin{frame}{大 $x$ 的主项}
  \[
    J_{\nu}(x)
    \sim\sqrt{\frac{2}{\pi x}}
      \cos\Bigl(x-\frac{\nu\pi}{2}-\frac{\pi}{4}\Bigr).
  \]
  振幅按 $x^{-1/2}$ 衰减，相邻零点的间距趋于 $\pi$。零点本身没有初等表达式。数值上 $J_0$ 的前两个正零点是
  \[
    j_{0,1}\approx 2.4048,
    \qquad
    j_{0,2}\approx 5.5201.
  \]
\end{frame}

\section{正交}

\begin{frame}{不同零点}
  设 $J_{\nu}(\alpha)=J_{\nu}(\beta)=0$，$\nu\geq 0$。$u(x)=J_{\nu}(\alpha x)$ 满足
  \[
    (xu')'+\Bigl(\alpha^2 x-\frac{\nu^2}{x}\Bigr)u=0,
  \]
  $v(x)=J_{\nu}(\beta x)$ 把 $\alpha$ 换成 $\beta$。相减后
  \[
    \frac{\mathrm{d}}{\mathrm{d}x}\bigl[x(uv'-vu')\bigr]
    =(\alpha^2-\beta^2)\,x\,uv.
  \]
  从 $0$ 积到 $1$。$x=0$ 处因子 $x$ 把边界项压成零，$x=1$ 处 $u=v=0$。$\alpha\neq\beta$ 时
  \[
    \int_0^{1} x\,J_{\nu}(\alpha x)\,J_{\nu}(\beta x)\,\mathrm{d}x=0.
  \]
  权重是 $x$。
\end{frame}

\begin{frame}{模}
  令 $\beta\to\alpha$。在 $J_{\nu}(\alpha)=0$ 处，$J_{\nu}'(\alpha)=-J_{\nu+1}(\alpha)$，
  \begin{equation*}
    \boxed{\displaystyle
      \int_0^{1} x\,J_{\nu}(\alpha x)^2\,\mathrm{d}x
      =\frac12 J_{\nu+1}(\alpha)^2.}
  \end{equation*}
  零点记成 $j_{\nu,m}$。区间改成 $[0,a]$、$J_{\nu}(ka)=0$ 时，模是
  \[
    \int_0^{a}\rho\,J_{\nu}(k\rho)^2\,\mathrm{d}\rho
    =\frac{a^2}{2}J_{\nu+1}(ka)^2.
  \]
\end{frame}

\begin{frame}{圆膜}
  半径 $a$、边缘固定、只看轴对称振动，
  \[
    u_{tt}
    =c^2\Bigl(u_{\rho\rho}+\frac{1}{\rho}u_{\rho}\Bigr),
    \qquad u(a,t)=0.
  \]
  分离后 $R=J_0(k\rho)$，$J_0(ka)=0$，所以 $k_m=j_{0,m}/a$，频率 $\omega_m=c k_m$。初速度为零时
  \[
    u(\rho,t)
    =\sum_{m=1}^{\infty} A_m J_0(k_m\rho)\cos(c k_m t),
  \]
  \[
    A_m
    =\frac{2}{a^2 J_1(k_m a)^2}
      \int_0^{a} f(\rho)\,J_0(k_m\rho)\,\rho\,\mathrm{d}\rho.
  \]
  基频是 $c j_{0,1}/a$。
\end{frame}

\begin{frame}{球贝塞尔函数}
  球坐标里的亥姆霍兹方程，径向因子是
  \[
    j_n(x)=\sqrt{\frac{\pi}{2x}}\,J_{n+1/2}(x).
  \]
  前两个是初等函数，
  \[
    j_0(x)=\frac{\sin x}{x},
    \qquad
    j_1(x)=\frac{\sin x}{x^2}-\frac{\cos x}{x}.
  \]
  $j_0(0)=1$。边界上 $j_0(ka)=0$ 就是 $\sin(ka)=0$，径向波数是 $n\pi/a$。
\end{frame}

\begin{frame}{这一讲留下的公式}
  \[
    J_{\nu}(x)
    =\sum_{k=0}^{\infty}
      \frac{(-1)^k}{k!\,\Gamma(k+\nu+1)}
      \Bigl(\frac{x}{2}\Bigr)^{2k+\nu}.
  \]
  \[
    J_{\nu-1}+J_{\nu+1}=\frac{2\nu}{x}J_{\nu},
    \qquad
    \int_0^{1} x J_{\nu}(\alpha x)^2\,\mathrm{d}x
    =\frac12 J_{\nu+1}(\alpha)^2.
  \]
  整数阶在原点的第二个解含对数。圆膜的径向函数是 $J_0(j_{0,m}\rho/a)$。
\end{frame}

\end{document}
